\subsubsection{问题二具体分析}

问题二在不使用储能的前提下，要求从"任务柔性格"与"电网交互机制"两个角度建立碳—成本最优的任务调度\cite{radovanovic2021carbon,goiri2015renewable}模型，且结论必须公平可比。本文将其抽象为：在约束（式 2--4）之上，引入购电、外送与弃电的电力平衡，把"任务迁移"直接映射为各区域逐时设施负荷的变化，从而进入统一的能源结算，并用一个可解释的身价信号（三段边际成本）引导自治优化，而非直接求解高维混合整数规划。

首先固定统一的结算口径（购电满足缺口、外送封顶、余量弃电），在问题一的可行基线上计算发出各区域的逐时边际成本面；接着按三段边际成本（弃电 0、外送 $\pi^s$、购电 $\pi^b$）对弹性任务做候选内重指派；然后冻结 $\lambda$ 面执行贪婪改进、反复重算 $\lambda$ 至收敛；最后以完整结算复算真实成本、碳与时延。该流程以边际成本面为桥梁，把任务柔性与电网交互统一到结算价值中。
